\chapter{结论}

\section{主要结论}

本文围绕 \code{gmsv} 源码，建立了一条从宠物成长、转生、融合、蛋喂药到战斗速度、命中、会心、属性状态和最终伤害的统一分析链。对宠物部分，本文明确区分了基础成长、波动后生效成长、蛋基础成长、孵化后成长与出生后随机状态，说明了 0 转宠与转生宠在 Rank 判定表上的本质差异，并把转生、融合、孵化与喂药都还原为带多次取整和上限裁切的离散映射。对战斗部分，本文整理了乱敏公式、骑宠混合公式、物理回避/命中/会心/反击公式、法术命中与伤害分层、场地与结界、属性加强、属性反转和若干职业状态变量的真实作用位置。

更重要的是，本文没有把这些机制停留在“规则说明”层，而是进一步给出了变量影响分析。命中与会心对敏捷并非全局单调；骑宠远程与骑宠近战的收益结构不同；结界强度在当前减伤实现里并未真正进入最终公式；属性反转必须区分“反转标记”与“当前固定属性快照”，不能简化成“连续命中两次立即复原”；蛋喂药问题则天然是一个离散优化问题，而不是简单的线性配比问题。这些结论共同说明，石器时代当前源码口径下的战斗与养成系统，本质上是一个由状态转移、分段函数和整数截断共同组成的复合系统。

\section{对玩家经验口径的修正}

本文得到的很多结论，都在不同程度上修正了玩家经验中的若干常见说法。其一，“敏捷高就一定更容易命中或更容易会心”并不成立，因为命中和会心对敏捷都存在局部非单调区间。其二，“结界强度越高减伤越多”在当前代码口径下并不严格成立，因为减伤分支主要读取的是结界是否存在，而不是存储下来的强度值。其三，“转生或融合只要总和更高就更强”同样不成立，因为 Rank 判表、取整顺序、单项上限和总和压缩都会显著改变最终收益。其四，“全部喂 5 级药一定最优”也不成立，因为压缩区中的整数份额分配常常会让低等级药在某些目标下更优。

这些修正并不是为了否定玩家经验，而是为了把经验重新放回源码约束之下。换言之，经验中有价值的部分通常对应某种局部近似，而源码分析的作用，是指出这套近似在哪些区间成立、又在哪些区间会失效。

\section{研究局限}

本文仍有若干明确的局限。首先，论文口径以当前持有的 \code{gmsv} 源码与项目实现为准，不同分支、不同编译宏或后续版本可能改变部分细节。其次，本文虽然建立了从养成到战斗的完整分析链，但在数值实验层面仍以结构分析和局部公式推导为主，尚未把所有章节都扩展成大规模模拟。再次，若干机制虽然已知主效果，但重复命中、跨状态覆盖与边界条件的全部分支还未被穷尽验证，尤其需要继续结合实战日志复核若干“源码直读尚不能充分解释”的现象。例如，白狼的连环攻击打在已被属性反转的目标上时，在未观察到防御动作减伤或偏转表现的前提下，两次伤害有时会出现极大差异；两只宠物使用同类连续命中技能时也可观察到类似现象。这类现象在体感上与“第二次受击时目标的反转属性未继续按第一次口径参与伤害结算”较为接近，但本文目前不把它写成既定结论，而只将其列为待进一步核对的研究问题。另一个同样需要保留为待研究问题的实战观察是：猎人装备弓箭并处于属性反转时，普通弓攻击与“剧毒武器”在伤害结算属性上可能表现出不同口径，且“剧毒武器”在多目标命中时还可能出现主目标与次目标表现不一致的情况。就当前静态源码来看，这两条路径虽然分属普通攻击分支与职业状态攻击分支，但最终都进入同一套物伤与属性修正链，尚不足以仅凭代码直读就严格解释上述差异，因此本文不把它写成已证实现象，而把它一并纳入后续待研究问题。

此外，本文为了保持可验证性，对一些当前源码中未实现或未能确认的机制都采取了保守表述。例如职业共通技能“自我强化”在当前代码中没有明确数值实现，因此本文没有替它补写一个假想公式。这种保守性会让部分章节看起来“结论偏少”，但它恰恰是源码研究的必要代价。

如果说本文的直接结果，是整理出一套较为严格的源码口径；那么它更长期的意义，则在于提供了一种方法：面对一个充满经验说法的老游戏系统，可以不依赖口耳相传，而是回到源码、变量、取整和状态转移本身，逐步建立出可验证、可复用、可计算的理论框架。这也是本文希望留下的最核心工作。

